\documentclass[english, parskip=full]{scrartcl}
\usepackage[utf8]{inputenc}  % Kodierung der Datei
\usepackage[ngerman]{babel}  % Sprache auf Deutsch 
\usepackage{color}
\usepackage{graphicx, gensymb}
\usepackage{amsfonts, cancel, amssymb}
\usepackage{amsmath, amsthm, array, esint}
\usepackage{booktabs, multirow}  % schönere Tabellen
\usepackage{caption}  % Anpassung der Captions
\usepackage{scrlayer-scrpage}    % Manipulation des Seitenstils
%\pagestyle{scrheadings}  % Stil für die Seite setzen
%\automark{section}  % Abschnittsnamen als Seitenbeschriftung 
%\setheadsepline{.5pt}    % Kopzeile durch Linie abtrennen
\usepackage[hidelinks]{hyperref}  % Links und weitere PDF-Features
\usepackage{typearea, tikz, pgffor, verbatim}
\usetikzlibrary{calc, quotes}
\usepackage[cache=false, outputdir=build]{minted}
\usepackage{placeins} % provides \FloatBarrier

\definecolor{bg}{rgb}{0.9,0.9,0.9}
\newcommand\numberthis{\addtocounter{equation}{1}\tag{\theequation}}
\newcommand{\bra}{}
\newcommand{\kom}{}
\newcommand{\brak}{}
\newcommand{\braket}{}
\newcommand{\intx}{}
\newcommand{\intr}{}
\newcommand{\kla}{}
\newcommand{\klam}{}
\newcommand{\klammer}{}
\newcommand{\oper}{}
\newcommand{\tecxt}{}
\newcommand{\Bra}{}
\newcommand{\Ket}{}
\renewcommand{\i}{\text{i}}
\newcommand{\e}{\text{e}}
\newcommand*{\Fourier}{\textsc{Fourier}\ }
\newcommand*{\F}{\mathcal{F}}
\newcommand*{\sinc}{\text{sinc\,}}
\newcommand{\conv}{\mathop{\scalebox{3.5}{\raisebox{-0.38ex}{\makebox[0.5\width][c]{$\ast$}}}}\limits}

\newcommand*{\titel}{03 Vorkonditionierung und Eigenwerte}
\newcommand*{\autor}{Joachim Schwardt und Julian Fleck}

\hypersetup{pdfauthor={\autor}, pdftitle={\titel}}

%\titlehead{titlehead}
\subject{Numerik II - Aufgaben}
\title{\titel}
%\author{\autor}
\author{\begin{tabular}{l|l}
Joachim Schwardt & 4768711 \\ 
Julian Fleck & 4759587\\ 
\end{tabular}}
\date{\today}

\begin{document}

%\begin{titlepage}
\maketitle  % Titel setzen
%\tableofcontents  % Inhaltsverzeichnis setzen
%\end{titlepage}

\section*{Aufgabe 1}
Gegeben sei eine symmetrische positiv definite matrix $A\in\mathbb{R}^{n\times n}$ und ein Vektor $b\in\mathbb{R}^n$.
Sei $D$ die Diagonale von $A$.
Die \textsc{Jacobi}-Vorkonditionierung verwendet $W:= D$. 
Seien die Diagonaleinträge nun alle gleich, also $D \equiv hI_{n\times n}$.
Dann lautet ein Iterationsschritt des vorkonditionierten CG-Verfahrens
\begin{align}
r_0 &= b - Ax_0 = \tilde{r}_0,\\
d_0 &= W^{-1}r_0 = \frac{1}{h}r_0 = \frac{1}{h} \tilde{d}_0,\\
\alpha_0 &= \frac{r_0^\top W^{-1} r_0}{d_0^\top Ad_0} = h\frac{r_0^\top r_0}{r_0^\top Ar_0} = h\tilde{\alpha}_0,\\
x_{1} &= x_0 + \alpha_0 d_0 = x_0 + \frac{r_0^\top r_0}{r_0^\top Ar_0} r_0 = x_0 + \tilde{\alpha}_0 \tilde{d}_0 = \tilde{x}_1,\\
r_1 &= r_0 - \alpha_0 Ad_0 = r_0 - \frac{r_0^\top r_0}{r_0^\top Ar_0} r_0 = \tilde{r}_0 - \tilde{\alpha}_0 \tilde{d}_0 = \tilde{r}_1,\\
\beta_1 &= \frac{r_1^\top r_1}{r_0^\top r_0} = \tilde{\beta}_1.
\end{align}
Die Variablen mit Tilde entsprechen dem Ergebnis ohne Vorkonditionierung. 
Im ersten Schritt ändert sich also nichts, und wir vermuten nun, dass obige Zusammenhänge allgemein gelten.
Die Induktionsvoraussetzung lautet also
\begin{align}
\beta_{k} &= \tilde{\beta}_{k},\\
r_k &= \tilde{r}_k,\\
d_k &= \frac{1}{h} \tilde{d}_k,\\
\alpha_k &= h\tilde{\alpha}_k,\\
x_{k+1} &= \tilde{x}_{k+1}.
\end{align}
Dann lautet ein Iterationsschritt des vorkonditionierten CG-Verfahrens
\begin{align}
r_{k+1} &= r_{k} - \alpha_kAd_k = \tilde{r}_k - h\tilde{\alpha}_k A \frac{1}{h} \tilde{d}_k = \tilde{r}_{k+1},\\
\beta_{k+1} &= \frac{r_{k+1}^\top r_{k+1}}{r_k^\top r_k} = \tilde{\beta}_{k+1},\\
d_{k+1} &= W^{-1}r_{k+1} + \beta_{k+1}d_{k} = \frac{1}{h}\tilde{r}_{k+1} + \tilde{\beta}_{k+1} \frac{1}{h} \tilde{d}_k = \frac{1}{h} \tilde{d}_{k+1},\\
\alpha_{k+1} &= \frac{r_{k+1}^\top W^{-1} r_{k+1}}{d_{k+1}^\top Ad_{k+1}} = h\frac{\tilde{r}_{k+1}^\top \tilde{r}_{k+1}}{\tilde{d}_k^\top A\tilde{d}_k} = h\tilde{\alpha}_{k+1},\\
x_{k+2} &= x_{k+1} + \alpha_{k+1} d_{k+1} = \tilde{x}_{k+1} + h \tilde{\alpha}_{k+1} \frac{1}{h} \tilde{d}_{k+1} = \tilde{x}_{k+2}.
\end{align}
Die Vorkonditionierung hat in diesem Fall also keinen Einfluss auf das CG-Verfahren.
(Das war irgendwie ziemlich umständlich, geht vermutlich deutlich kürzer...)

\section*{Aufgabe 2}
Wir betrachten einen Algorithmus zur Berechnung der LR-Zerlegung einer Matrix $A\in\mathbb{R}^{n\times n}$, wobei $L$ eine linke untere Dreiecksmatrix mit $L_{jj} = 1$ und $R$ eine rechte obere Dreiecksmatrix sein sollen, sodass $LR=A$.
%(Anmerkung: Wir glauben dass die innerste Summe von $j=1$ bis $j=k$ (statt $k-1$) bzw. $j=1$ bis $j=i$ (statt $i-1$) laufen sollte, weil bei uns sonst der Algorithmus nicht funktioniert.)
%Wir wollen zunächst zeigen, dass der Algorithmus für symmetrische positiv definite Matrizen immer explizit durchführbar ist.
%Wir wollen zunächst zeigen, dass für beliebige Matrizen immer alle Informationen vorliegen, um den Algorithmus (prinzipiell) durchzuführen.
%Hier ignorieren wir noch, ob durch Null geteilt wird.
%Wir unterschiden also nur zwischen besetzten ($\ast$) und unbesetzten (0) Einträgen.
%Neue Einträge sind blau markiert, und Einträge, die im ächsten Schritt benötigt werden, sind unterstrichen.
%\\
%Initialisierung:
%\begin{align}
%L &= \begin{pmatrix}
%\color{red}{\ast} & 0 & \dots & 0 \\
%0 & \ast & \dots & 0 \\
%\vdots & \vdots & \ddots & \vdots \\
%0 & 0 & \dots & \ast
%\end{pmatrix} &&\tecxt\text{und} & R &= \begin{pmatrix}
%0 & 0 & \dots & 0  \\
%0 & 0 & \dots & 0 \\
%\vdots & \vdots & \ddots & \vdots \\
%0 & 0 & \dots & 0
%\end{pmatrix}.
%\end{align}
%Schritt $i=1$: $k=1,\dots,0 \rightarrow$ entfällt. 
%Für $k=1,\dots,n$:
%\begin{align}
%R_{1k} &= A_{1k} - \sum_{j=1}^0 L_{1j}R_{jk} = A_{1k}
%\end{align}
%Jetzt ist
%\begin{align}
%L &= \begin{pmatrix}
%\ast & 0 & \dots & 0 \\
%0 & \ast & \dots & 0 \\
%\vdots & \vdots & \ddots & \vdots \\
%0 & 0 & \dots & \ast
%\end{pmatrix} &&\tecxt\text{und} & R &= \begin{pmatrix}
%\underline{\color{blue}{\ast}} & \color{blue}{\ast} & \dots & \color{blue}{\ast}  \\
%0 & 0 & \dots & 0 \\
%\vdots & \vdots & \ddots & \vdots \\
%0 & 0 & \dots & 0
%\end{pmatrix}.
%\end{align}
%Schritt $i=2$: 
%Für $k=1$:
%\begin{align}
%L_{2k} &= \frac{1}{R_{kk}} \kla\left( A_{2k} - \sum_{j=1}^{k-1}L_{2j}R_{jk} \right)
%\end{align} 
%Jetzt ist
%\begin{align}
%L &= \begin{pmatrix}
%\ast & 0 & \dots & 0 \\
%\underline{\color{blue}{\ast}} & \ast & \dots & 0 \\
%\vdots & \vdots & \ddots & \vdots \\
%0 & 0 & \dots & \ast
%\end{pmatrix} &&\tecxt\text{und} & R &= \begin{pmatrix}
%\underline{\ast} & \ast & \dots & \ast  \\
%0 & 0 & \dots & 0 \\
%\vdots & \vdots & \ddots & \vdots \\
%0 & 0 & \dots & 0
%\end{pmatrix}.
%\end{align}
%und für $k=2,\dots,n$:
%\begin{align}
%R_{2k} &= A_{2k} - \sum_{j=1}^1 L_{2j}R_{jk}
%\end{align}
%Jetzt ist
%\begin{align}
%L &= \begin{pmatrix}
%\ast & 0 & \dots & 0 \\
%\ast & \ast & \dots & 0 \\
%\vdots & \vdots & \ddots & \vdots \\
%0 & 0 & \dots & \ast
%\end{pmatrix} &&\tecxt\text{und} & R &= \begin{pmatrix}
%\ast & \ast & \dots & \ast  \\
%0 & \color{blue}{\ast} & \dots & \color{blue}{\ast} \\
%\vdots & \vdots & \ddots & \vdots \\
%0 & 0 & \dots & 0
%\end{pmatrix}.
%\end{align}
%Schritt $i=3$: 
%Für $k=1$:
%\begin{align}
%L_{3k} &= \frac{1}{R_{kk}} \kla\left( A_{3k} - \sum_{j=1}^{k-1}L_{3j}R_{jk} \right)
%\end{align} 
%Jetzt ist
%\begin{align}
%L &= \begin{pmatrix}
%\ast & 0 & \dots & 0 \\
%\ast & \ast & \dots & 0 \\
%\vdots & \vdots & \ddots & \vdots \\
%0 & 0 & \dots & \ast
%\end{pmatrix} &&\tecxt\text{und} & R &= \begin{pmatrix}
%\ast & \ast & \dots & \ast  \\
%0 & \color{blue}{\ast} & \dots & \color{blue}{\ast} \\
%\vdots & \vdots & \ddots & \vdots \\
%0 & 0 & \dots & 0
%\end{pmatrix}.
%\end{align}
%Für $k=2$:
%\begin{align}
%L_{3k} &= \frac{1}{R_{kk}} \kla\left( A_{3k} - \sum_{j=1}^{k-1}L_{3j}R_{jk} \right)
%\end{align} 
%Jetzt ist
%\begin{align}
%L &= \begin{pmatrix}
%\ast & 0 & \dots & 0 \\
%\underline{\color{blue}{\ast}} & \ast & \dots & 0 \\
%\vdots & \vdots & \ddots & \vdots \\
%0 & 0 & \dots & \ast
%\end{pmatrix} &&\tecxt\text{und} & R &= \begin{pmatrix}
%\underline{\ast} & \ast & \dots & \ast  \\
%0 & 0 & \dots & 0 \\
%\vdots & \vdots & \ddots & \vdots \\
%0 & 0 & \dots & 0
%\end{pmatrix}.
%\end{align}
%und für $k=3,\dots,n$:
%\begin{align}
%R_{3k} &= A_{3k} - \sum_{j=1}^1 L_{3j}R_{jk}
%\end{align}
%Jetzt ist
%\begin{align}
%L &= \begin{pmatrix}
%\ast & 0 & \dots & 0 \\
%\ast & \ast & \dots & 0 \\
%\vdots & \vdots & \ddots & \vdots \\
%0 & 0 & \dots & \ast
%\end{pmatrix} &&\tecxt\text{und} & R &= \begin{pmatrix}
%\ast & \ast & \dots & \ast  \\
%0 & \color{blue}{\ast} & \dots & \color{blue}{\ast} \\
%\vdots & \vdots & \ddots & \vdots \\
%0 & 0 & \dots & 0
%\end{pmatrix}.
%\end{align}
\subsection*{a) Durchführbarkeit des Algorithmus}
Für $i=1$ wird nur die erste Zeile von $A$ in $R$ kopiert, dieser Schritt ist also immer durchführbar.
Angenommen wir haben die ersten $i$ Schritte durchgeführt.
Jetzt haben die Matrizen also die Form
\begin{align}
L &= \begin{pmatrix}
1 & 0 & \dots & 0 & 0 & \dots & 0 \\
\ast & 1 & \dots & 0 & 0 & \dots & 0 \\
\vdots & \vdots & \ddots & \vdots & \vdots & \ddots & \vdots \\
\ast_{i1} & \ast_{i2} & \dots & 1_{ii} & 0 & \dots & 0 \\
0 & 0 & \dots & 0 & 1 & \dots & 0 \\
\vdots & \vdots & \vdots & \vdots & \vdots & \ddots & \vdots \\
0 & 0 & \dots & 0 & 0 & \dots & 1
\end{pmatrix}, && & R &= \begin{pmatrix}
\ast & \ast & \dots & \ast & \ast & \dots & \ast \\
0 & \ast & \dots & \ast & \ast & \dots & \ast \\
\vdots & \vdots & \ddots & \vdots & \vdots & \ddots & \vdots \\
0 & 0 & \dots & \ast_{ii} & \ast_{i,i+1} & \dots & \ast_{in} \\
0 & 0 & \dots & 0 & 0 & \dots & 0 \\
\vdots & \vdots & \vdots & \vdots & \vdots & \ddots & \vdots \\
0 & 0 & \dots & 0 & 0 & \dots & 0
\end{pmatrix}.
\end{align}
Dann betrachten wir Schritt $i+1$, wobei zunächst $k=1,\dots,i$:
\begin{align}
L_{i+1,k} &= \frac{1}{R_{kk}} \kla\left( A_{ik} - \sum_{j=1}^{k-1} L_{i+1,j}R_{jk} \right).
\end{align}
Wir benötigen hier die ersten $i$ Spalten von $R$, und davon jeweils die Zeilen $j=1,\dots,k-1$.
Insbesondere ist der Zeilenindex also immer kleiner als der Spaltenindex, womit alle Informationen bezüglich $R$ hier vorhanden sind.
Von $L$ benötigen wir im Schritt $k$ (also für $L_{i+1,k}$) immer die ersten $j=1,\dots,k-1$ Spalten der Zeile $i+1$.
Damit sind also auch hier immer die Einträge vorhanden (es wird in jedem Schritt zuletzt der Eintrag $L_{i+1,k-1}$ verwendet, welcher gerade im Schritt davor berechnet wird).
\\
Im zweiten Teil des Algorithmus ist $k=i+1,\dots,n$:
\begin{align}
R_{i+1,k} &= A_{i+1,k} - \sum_{j=1}^{i} L_{i+1,j}R_{jk}.
\end{align}
Wir benötigen hier von $R$ die Spalten $k=i+1,\dots,n$, und davon jeweils die Zeilen $j=1,\dots,i$.
Auch hier ist $k>j$, es sind also alle relevanten Einträge bekannt.
Von $L$ benötigen wir in der Zeile $i+1$ die Einträge $j=1,\dots,i$, welche alle im ersten Teil von Schritt $i+1$ des Algorithmus bestimmt wurden.
Damit sollte gezeigt sein, dass die notwendigen Informationen immer vorliegen.
\\
Für die Durchführbarkeit fehlt aber noch $R_{kk} \neq 0$, was insbesondere mit der positiven Definitheit zusammenhängen muss.
Das kann nur im zweiten Teil passieren, und dort nur für $k=i+1$.
Hier gilt
\begin{align}
R_{i+1,i+1} &= A_{i+1,i+1} - \sum_{j=1}^i L_{i+1,j}R_{j,i+1}.
\end{align}
Jetzt müsste man vermutlich irgendetwas sehen... 
\\ 
Wir wissen, dass $A_{jj} > 0$ für pos. def. Matrizen, aber das reicht natürlich noch nicht für $R_{kk}\neq 0$.
Da $A$ s.p.d., können wir $R=DL^\top$ schreiben, wobei $D$ eine Diagonalmatrix ist.
Dann erhalten wir (unter Verwendung von $L_{jj} = 1$) den Zusammenhang
\begin{align}
D_{i+1,i+1} &= A_{i+1,i+1} - \sum_{j=1}^i L_{i+1,j}^2 D_{jj}.
\end{align}
(Jetzt müsste $D_{i+1,i+1} \neq 0$ sein, damit kommen wir aber auch nicht wirklich weiter.)
%Die andere Gleichung ergibt für $k=j$ dann
%\begin{align}
%L_{ij} &= \frac{1}{D_{jj}} \kla\left( A_{ij} - \sum_{l=1}^{j-1} L_{il} D_{ll} L_{jl} \right)•
%\end{align}
%Andererseits ist $A=LDL^\top$, also $A_{jj} = \sum_{k=1}^n L_{jk}D_{kk}L_{jk}$.
%Damit folgt
%\begin{align}
%D_{i+1,i+1} &= \sum_{j=1}^n L_{i+1,j}^2 D_{jj} - \sum_{j=1}^i L_{i+1,j}^2 D_{jj} = \sum_{j=i+1}^n L_{i+1,j}^2 D_{jj} = L_{i+1,i+1}D_{i+1,i+1}
%\end{align}

\subsection*{b) Spezialfall}
Für den speziellen Fall
\begin{align}
A &= \begin{pmatrix}
1 & 2 & 5 \\ 2 & 4 & 2 \\ 5 & 2 & 2
\end{pmatrix}
\end{align}
schlägt der Algorithmus fehl. 
Die Matrix ist zwar symmetrisch und invertierbar, da
\begin{align}
\det A &= 8 + 20 + 20 - 100 - 4 - 8 \neq 0,
\end{align}
aber sie ist nicht positiv definit.
Initialisierung:
\begin{align}
L &= \begin{pmatrix}
1 & 0 & 0 \\ 0 & 1 & 0 \\ 0 & 0 & 1
\end{pmatrix} && \tecxt\text{und} & R &= \begin{pmatrix}
0 & 0 & 0 \\ 0 & 0 & 0 \\ 0 & 0 & 0
\end{pmatrix}.
\end{align}
Schritt $i=1$: $k=1,\dots,0 \rightarrow$ entfällt. 
Für $k=1,\dots,3$:
\begin{align}
R_{1k} &= A_{1k} - \sum_{j=1}^0 L_{1j}R_{jk} = A_{1k}
\end{align}
Jetzt ist
\begin{align}
L &= \begin{pmatrix}
1 & 0 & 0 \\ 0 & 1 & 0 \\ 0 & 0 & 1
\end{pmatrix} && \tecxt\text{und} & R &= \begin{pmatrix}
1 & 2 & 5 \\ 0 & 0 & 0 \\ 0 & 0 & 0
\end{pmatrix}.
\end{align}
Schritt $i=2$: Für $k=1,\dots,1$:
\begin{align}
L_{21} &= \frac{1}{R_{11}} \kla\left( A_{21} - \sum_{j=1}^{0}L_{2j}R_{j1} \right) = 2
\end{align} 
und für $k=2,\dots,3$:
\begin{align}
R_{22} &= A_{22} - \sum_{j=1}^1 L_{2j}R_{j2} = 4 - 2\cdot 2 = 0,\\
R_{23} &= \dots
\end{align}
Da $R_{22} = 0$ wird im nächsten Schritt durch Null geteilt, womit der Algorithmus fehlschlägt.

%A2 = np.array([[1,2,5],[2,4,2],[5,2,2]])    # NICHT POSITIV DEFINIT!!
%
%def lr_decomp(A):
%    n = A.shape[0]
%    L = np.eye(n)
%    R = np.zeros((n,n))
%    for i in range(n):
%        print(f"{i = }")
%        for k in range(i):
%            print("INNER FIRST:", [(L[i, j], R[j, k]) for j in range(k)])
%            L[i, k] = (A[i, k] - np.sum([L[i, j] * R[j, k]  for j in range(k)])) / R[k, k]
%            print(f"FIRST: Write to L[{i},{k}] = {L[i, k]}")
%        for k in range(i, n):
%            print("INNER SECOND:", [(L[i, j], R[j, k]) for j in range(i)])
%            R[i, k] = A[i, k] - np.sum([L[i, j] * R[j, k] for j in range(i)])
%            print(f"SECOND: Write to R[{i},{k}]= {R[i, k]}")
%        print()
%    return L, R
%
%L, R = lr_decomp(A2)
%print(L, R, L@R)


\section*{Aufgabe 3}
Sei $A\in\mathbb{C}^{n\times n}$.
Sei $A^\dagger := (A^\ast)^\top$, wobei ${}^\ast$ die elementweise komplexe Konjugation beschreibt.
Sei $\lambda\in\mathbb{C}$ ein Eigenwert von $A$.
Dann gilt
\begin{align}
0 &= \det \kla\left( A - \lambda I_{n\times n} \right) = \chi_A(\lambda).
\end{align}
Da $\chi_A$ ein Polynom in $\lambda$ ist, bei dem die Koeffizienten von den Einträgen von $A$ abhängen, gilt 
\begin{align}
\chi_A(\lambda) &= \sum_{k=0}^n \lambda^k c_k(A).
\end{align}
Dabei enthalten die $c_k$ nur Summen und Produkte der $A_{jk}$, weshalb $c_k(A)^\ast = c_k(A^\ast)$ gilt.
Daraus folgt dann
\begin{align}
\chi_A(\lambda)^\ast &= \sum_{k=0}^n \kla\left( \lambda^\ast \right)^k c_k(A^\ast) = \chi_{A^\ast}(\lambda^\ast).
\end{align}
Als Determinante geschrieben ergibt sich dann
\begin{align}
0 &= \det\kla\left( A^\ast - \lambda^\ast I_{n\times n} \right) \kom\overset{\substack{{\det B = \det B^\top} \\ |}}{=} \det\kla\left( A^\dagger - \lambda^\ast I_{n\times n} \right),
\end{align}
womit also $\lambda^\ast$ ein Eigenwert von $A^\dagger$ ist.

\section*{Aufgabe 4}
Gegeben ist die Matrix
\begin{align}
B &= \begin{pmatrix}
3 & 0 & 0 \\ -1 & 1 & 2 \\ 1 & -2 & 1
\end{pmatrix}.
\end{align}
Das charakteristische Polynom lautet
\begin{align}
\chi(\lambda) &= \det\kla\left( B - \lambda I_{3\times 3} \right) = \det\begin{pmatrix}
3-\lambda & 0 & 0 \\ -1 & 1-\lambda & 2 \\ 1 & -2 & 1-\lambda
\end{pmatrix} = \\
&= (3-\lambda) \det\begin{pmatrix}
1-\lambda & 2 \\ -2 & 1 -\lambda
\end{pmatrix} = \\
&= (3-\lambda) \kla\left( (1-\lambda)^2 + 4 \right) \overset{!}{=} 0.
\end{align}
Ein Eigenwert lässt sich also direkt als $\lambda_1=3$ ablesen.
Die anderen beiden ergeben sich aus $0 = \lambda^2 - 2\lambda + 5$, woraus 
\begin{align}
\lambda_{2,3} = 1 \pm \sqrt{1 - 5} = 1 \pm 2\i,
\end{align}
folgt.
Also sind $\lambda_2 = 1 + 2\i$ und $\lambda_3 = 1 - 2\i$.
Sei $v=(a,b,c)^\top$ ein Eigenvektor. 
Dann gilt für $\lambda=\lambda_1$
\begin{align}
0 &= \begin{pmatrix}
3-\lambda & 0 & 0 \\ -1 & 1-\lambda & 2 \\ 1 & -2 & 1-\lambda
\end{pmatrix} \begin{pmatrix}
a \\ b \\ c
\end{pmatrix} = \begin{pmatrix}
0 & 0 & 0 \\ -1 & -2 & 2 \\ 1 & -2 & -2
\end{pmatrix} \begin{pmatrix}
a \\ b \\ c
\end{pmatrix},
\end{align}
also $a=2(c-b)$ und $a=2(c+b)$.
Daraus folgt $b=0$ und $a=2c$.
Für $\lambda=\lambda_{2,3}$
\begin{align}
0 &= \begin{pmatrix}
3-\lambda & 0 & 0 \\ -1 & 1-\lambda & 2 \\ 1 & -2 & 1-\lambda
\end{pmatrix} \begin{pmatrix}
a \\ b \\ c
\end{pmatrix} = \begin{pmatrix}
2 \mp 2\i & 0 & 0 \\ -1 & \mp 2\i & 2 \\ 1 & -2 & \mp 2\i
\end{pmatrix} \begin{pmatrix}
a \\ b \\ c
\end{pmatrix},
\end{align} 
also $a=0$ und $b=\mp \i c$.
Die Eigenvektoren lauten dann
%\begin{align}
%v_1 &= \frac{1}{\sqrt{5}}\begin{pmatrix}
%2 \\ 0 \\ 1
%\end{pmatrix}, &&& v_2 &= \frac{1}{\sqrt{2}}\begin{pmatrix}
%0 \\ -\i \\ 1
%\end{pmatrix}, &&& v_3 &= \frac{1}{\sqrt{2}}\begin{pmatrix}
%0 \\ \i \\ 1
%\end{pmatrix}.
%\end{align}
\begin{align}
v_1 &= \begin{pmatrix}
2 \\ 0 \\ 1
\end{pmatrix}, &&& v_2 &= \begin{pmatrix}
0 \\ -\i \\ 1
\end{pmatrix}, &&& v_3 &= \begin{pmatrix}
0 \\ \i \\ 1
\end{pmatrix}.
\end{align}
Allerdings ist $v_1$ nicht orthogonal zu $v_{2,3}$.
Die Linearkombination $v_1 - \frac{1}{2}(v_2+v_3) = (2,0,0)^\top$ erfüllt dies, womit wir eine orthonormale Basis gefunden haben.
Diese lautet
\begin{align}
a_1 &= \begin{pmatrix}
1 \\ 0 \\ 0
\end{pmatrix}, &&& a_2 &= \frac{1}{\sqrt{2}}\begin{pmatrix}
0 \\ -\i \\ 1
\end{pmatrix}, &&& a_3 &= \frac{1}{\sqrt{2}}\begin{pmatrix}
0 \\ \i \\ 1
\end{pmatrix}.
\end{align}
%Die Transformationsmatrix können wir dann als
%\begin{align}
%S &= \frac{1}{\sqrt{2}} \begin{pmatrix}
%\sqrt{2} & 0 & 0 \\ 0 & -\i & \i \\ 0 & 1 & 1
%\end{pmatrix}
%\end{align}
%schreiben, wobei $S^{-1} = S^\dagger$.
%Es gilt $B = S DS^\dagger$, wobei $D=\tecxt\text{diag\,}(\lambda_1,\lambda_2,\lambda_3)$
%\begin{align}
%S DS^\dagger &= \frac{1}{2} \begin{pmatrix}
%\sqrt{2} & 0 & 0 \\ 0 & -\i & \i \\ 0 & 1 & 1
%\end{pmatrix}\begin{pmatrix}
%3 & 0 & 0 \\ 0 & 1+2\i & 0 \\ 0 & 0 & 1-2\i
%\end{pmatrix}\begin{pmatrix}
%\sqrt{2} & 0 & 0 \\ 0 & \i & 1 \\ 0 & -\i & 1
%\end{pmatrix} = \\
%&= \frac{1}{2} \begin{pmatrix}
%\sqrt{2} & 0 & 0 \\ 0 & -\i & \i \\ 0 & 1 & 1
%\end{pmatrix}\begin{pmatrix}
%3\sqrt{2} & 0 & 0 \\ 0 & -2+\i & 1+2\i \\ 0 & -2-\i & 1-2\i
%\end{pmatrix} = \\
%&= \frac{1}{2} \begin{pmatrix}
%6 & 0 & 0 \\ 0 & 2 & 4 \\ 0 & -4 & 2
%\end{pmatrix} = \begin{pmatrix}
%3 & 0 & 0 \\ 0 & 1 & 2 \\ 0 & -2 & 1
%\end{pmatrix} \neq B.
%\end{align}
Für die Transformationsmatrix können wir die $a_{1,2,3}$ allerdings nicht direkt verwenden.
Stattdessen ist
\begin{align}
S &= \begin{pmatrix}
2 & 0 & 0 \\ 0 & -\i & \i \\ 1 & 1 & 1
\end{pmatrix}.
\end{align}
Wir berechnen die Inverse
\begin{align}
S= \begin{pmatrix}
2 & 0 & 0 \\ 0 & -\i & \i \\ 1 & 1 & 1
\end{pmatrix} &\Bigg\vert \begin{pmatrix}
1 & 0 & 0 \\ 0 & 1 & 0 \\ 0 & 0 & 1
\end{pmatrix},\\
\begin{pmatrix}
1 & 0 & 0 \\ 0 & 1 & -1 \\ 0 & 1 & 1
\end{pmatrix} &\Bigg\vert \begin{pmatrix}
\frac{1}{2} & 0 & 0 \\ 0 & \i & 0 \\ -\frac{1}{2} & 0 & 1
\end{pmatrix},\\
\begin{pmatrix}
1 & 0 & 0 \\ 0 & 1 & -1 \\ 0 & 0 & 1
\end{pmatrix} &\Bigg\vert \begin{pmatrix}
\frac{1}{2} & 0 & 0 \\ 0 & \i & 0 \\ -\frac{1}{4} & -\frac{\i}{2} & \frac{1}{2}
\end{pmatrix},\\
\begin{pmatrix}
1 & 0 & 0 \\ 0 & 1 & 0 \\ 0 & 0 & 1
\end{pmatrix} &\Bigg\vert \begin{pmatrix}
\frac{1}{2} & 0 & 0 \\ -\frac{1}{4} & \frac{\i}{2} & \frac{1}{2} \\ -\frac{1}{4} & -\frac{\i}{2} & \frac{1}{2}
\end{pmatrix} = S^{-1}.
\end{align}
Es gilt $B = S DS^{-1}$, wobei $D=\tecxt\text{diag\,}(\lambda_1,\lambda_2,\lambda_3)$
\begin{align}
S DS^{-1} &= \begin{pmatrix}
2 & 0 & 0 \\ 0 & -\i & \i \\ 1 & 1 & 1
\end{pmatrix}\begin{pmatrix}
3 & 0 & 0 \\ 0 & 1+2\i & 0 \\ 0 & 0 & 1-2\i
\end{pmatrix}\frac{1}{2} \begin{pmatrix}
1 & 0 & 0 \\ -\frac{1}{2} & \i & 1 \\ -\frac{1}{2} & -\i & 1
\end{pmatrix} = \\
&= \frac{1}{2} \begin{pmatrix}
2 & 0 & 0 \\ 0 & -\i & \i \\ 1 & 1 & 1
\end{pmatrix}\begin{pmatrix}
3 & 0 & 0 \\ -\i-\frac{1}{2} & \i-2 & 1+2\i \\ \i - \frac{1}{2} & -\i-2 & 1 - 2\i
\end{pmatrix} = \\
&= \frac{1}{2} \begin{pmatrix}
6 & 0 & 0 \\ -2 & 2 & 4 \\ 2 & -4 & 2
\end{pmatrix} = B.
\end{align}
Daraus folgt insbesondere auch $B^k = (SDS^{-1})^k = SDS^{-1}SDS^{-1} \dots SDS^{-1} = SD^kS^{-1}$.
Damit lässt sich das Matrixexponential dann einfach berechnen, da
\begin{align}
\exp(B) &= \sum_{k\ge 0} \frac{1}{k!}B^k = S \kla\left( \sum_{k\ge 0} \frac{1}{k!}D^k \right) S^{-1} = S \begin{pmatrix}
\e^{\lambda_1} & 0 & 0 \\
0 & \e^{\lambda_2} & 0 \\
0 & 0 & \e^{\lambda_3}
\end{pmatrix} S^{-1} = \\
&= \begin{pmatrix}
2 & 0 & 0 \\ 0 & -\i & \i \\ 1 & 1 & 1
\end{pmatrix}\begin{pmatrix}
\e^3 & 0 & 0 \\ 0 & \e^{1+2\i} & 0 \\ 0 & 0 & \e^{1-2\i}
\end{pmatrix} \frac{1}{2} \begin{pmatrix}
1 & 0 & 0 \\ -\frac{1}{2} & \i & 1 \\ -\frac{1}{2} & -\i & 1
\end{pmatrix} = \\
&= \frac{1}{2} \begin{pmatrix}
2 & 0 & 0 \\ 0 & -\i & \i \\ 1 & 1 & 1
\end{pmatrix} \begin{pmatrix}
\e^3 & 0 & 0 \\ 
-\frac{\e}{2} \e^{2\i} & \e\i\e^{2\i} & \e\e^{2\i} \\
-\frac{\e}{2} \e^{-2\i} & -\e\i\e^{-2\i} & \e\e^{-2\i}
\end{pmatrix} = \\
&= \frac{1}{2} \begin{pmatrix}
2\e^3 & 0 & 0 \\
\frac{\e\i}{2}\kla\left( \e^{2\i} - \e^{-2\i} \right) & \e\kla\left( \e^{2\i} + \e^{-2\i} \right) & \e\i\kla\left( -\e^{2\i} + \e^{-2\i} \right) \\
\e^3 - \frac{\e}{2}\kla\left( \e^{2\i} + \e^{-2\i} \right) & \e\i \kla\left( \e^{2\i} - \e^{-2\i} \right) & \e\kla\left( \e^{2\i} + \e^{-2\i} \right)
\end{pmatrix} = \\
&= \begin{pmatrix}
\e^3 & 0 & 0 \\
-\e\sin(2) & 2\e\cos(2) & 2\e\sin(2) \\
\frac{1}{2}\e^3 - \frac{1}{2}\e\cos(2) & -2\e\sin(2) & 2\e\cos(2)
\end{pmatrix}.
\end{align}
Im Allgemeinen ist $S$ nicht dünnbesetzt, da die Eigenvektoren einer nicht notwendiger dünnbesetzt sind.
Daher eignet sich dieses Vorgehen schlecht für große dünnbesetzte Matrizen $B$.

%B = np.array([[3,0,0],[-1,1,2],[1,-2,1]])
% scipy.linalg.expm(B)
% expB = np.array([[np.exp(3), 0, 0], [-np.e * np.sin(2), 2 * np.e * np.cos(2), 2*np.e*np.sin(2)], [0.5 * np.exp(3) - 0.5*np.e * np.cos(2), -2*np.e*np.sin(2), 2*np.e*np.cos(2)]]) #### FALSCH!!!
%S = np.array([[np.sqrt(2),0,0],[0,-1j,1j],[0,1,1]]) / np.sqrt(2)


\section*{Programmieraufgabe 3 (schriftlicher Teil)}
Bei der unvollständigen LR-Zerlegungn lassen wir im Algorithmus alle Einträge weg, bei denen $A_{ik} = 0$.
Insbesondere werden alle Summen nur über solche $L_{ij},R_{ij}$ ausgeführt, für die $A_{ij}\neq 0$.
Im Folgenden lassen wir die Tilde über den Matrizen weg.
\\
Mit dem Zusammenhang $R=DL^\top$ folgt $R_{jk} = D_{jj}L_{kj}$.
Einsetzen in den Algorithmus liefert dann unter Ausnutzen von $L_{jj}=1$
\begin{align}
L_{ik} &= \frac{1}{R_{kk}} \kla\left( A_{ik} - \sum_{j=1}^{k-1}L_{ij}R_{jk} \right) = \frac{1}{D_{kk}} \kla\left( A_{ik} - \sum_{j=1}^{k-1}L_{ij}D_{jj}L_{kj} \right).
\end{align}
Der zweite Teil lautet unter Verwendung von $A=A^\top$
\begin{align}
R_{ik} &= A_{ik} - \sum_{j=1}^{i-1} L_{ij}R_{jk},\\
\Rightarrow\ \ \ D_{ii}L_{ki} &= A_{ki} - \sum_{j=1}^{i-1} L_{ij}D_{jj}L_{kj}, \label{eqn:R ik incomplete choelsky} \\
\Rightarrow\ \ \ L_{i'k'} &= \frac{1}{D_{k'k'}} \kla\left( A_{i'k'} - \sum_{j=1}^{k'-1} L_{k'j}D_{jj}L_{i'j} \right). \label{eqn:L ik incomplete choelsky} 
\end{align} 
Hier haben wir die Indizes umbenannt ($i\rightarrow k'$ und $k\rightarrow i'$).
In dieser Form entspricht \autoref{eqn:L ik incomplete choelsky} genau dem vorherigen Ergebnis, es müssen ledigleich die Striche der Indizes weggelassen und die Reihenfolge der $L$-Terme in der Summe vertauscht werden.
Inbesondere kann dieser Teil also einfach weggelassen werden, da er bereits im ersten Teil berechnet wird.
Allerdings fehlt uns dann noch der Spezialfall von $k=i$.
Dann gilt $R_{ii} = D_{ii}L_{ii} = D_{ii}$ und wir können direkt \autoref{eqn:R ik incomplete choelsky} verwenden, also
\begin{align}
D_{ii} &= A_{ii} - \sum_{j=1}^{i-1}L_{ij}D_{jj}L_{ij}.
\end{align}

%\begin{thebibliography}{9}
%\bibitem{example} description, \url{https://www.exmapleurl}, day.\,month~2021.
%\end{thebibliography}

\end{document}
